\subsection{The two tensor factors and their integral flags}
\label{sec:count-three-tensor}

With the base field fixed, the four residual labels prescribe ratios for the two tensor factors. The integral structures on those factors will determine which extensions occur on reduction.

Use \(F_0,K,K',g_0\) from Section~\ref{sec:count-three-tower}.
Write \(\sigma=1\) for the original order and \(\sigma=-1\)
for the exchanged order. Let \(D_i\) denote the ordered
characteristic-zero characters of the diagonal rational datum,
and let \(d_i\) be their reductions. Their quotient
\(d_0/d_1\) on \(G_{F_0}\) is \(\lambda_0\).
This is the character of \(K'/\Q\), and it is nontrivial
at \(g_0\). In particular \(K'\) is contained in neither
\(F_0\) nor \(F_0K\), and both exclusions persist on the
\(\ell\)-power layers by taking an appropriate \(\ell\)-power of \(g_0\).

\subsubsection{The quadratic factor}

The first factor must distinguish the two orders of the tensor labels. Inducing an algebraic Hecke character from the auxiliary quadratic field prescribes both its inertia weights and its action at the chosen normalizer.

\begin{proposition}[The quadratic integral factor]\label{prop:count-three-quadratic-factor}
The algebraic Hecke character of \(K'\) given by the
class-number construction yields the integral rank-two
quadratic factor, with the specified diagonal action at \(g_0\). Its two inertia exponents at the chosen place above \(\ell\) are 
\[
 (m,-m).
\]

\end{proposition}

\begin{proof}
Let \(h'\) be the class number of \(K'\). For each fractional
ideal \(\mathfrak a\), choose a generator \(a\) for
\(\mathfrak a^{h'}\). For a positive multiple \(r\) of
the number of roots of unity of \(K'\), the character
\[
                 \mathfrak a\longmapsto (a/\bar a)^r
\]
is independent of the generator because changing \(a\) by a
root of unity multiplies this value by an \(r\)-th power equal
to one. Multiplication of the ideals multiplies the generators
up to such a unit. On a principal ideal \((x)\) its value is
\[
                       (x/\bar x)^{rh'}.
\]
These exponents are its infinity type.

It is unramified at every finite prime as a Hecke character
and conjugation inverts it. Its \(\ell\)-adic avatar
\(\kappa'\) is unramified away from \(\ell\).
Above \(\ell\), the two places are split, and on their
inertia groups the exponents are \(m_0,-m_0\), where
\(m_0=rh'\), up to choosing the first place.
The avatar takes values in the finite local coefficient
field given by the chosen embedding of \(K'\).
Its values are units by continuity and compactness.
Raise this character to a positive power to make its
reduction one and keep \(m_0\) for the resulting
positive exponent. The inertia exponents are then
still opposite, and conjugation still inverts the
character.

For any positive integer \(a\), take
\[
 \tau=\operatorname{Ind}_{G_{K'}}^{G_{\Q}} (\kappa')^a
\]

and write \(m=am_0\). Since \(g_0\notin G_{K'}\), conjugation by it inverts
\(\kappa'\) and fixes \(g_0^2\). Thus
\[
 \kappa'(g_0^2)=\kappa'(g_0^2)^{-1}.
\]
This value reduces
to one, so it is one because \(\ell\) is odd.
Consequently in the two induced basis vectors the action
of \(g_0\) interchanges the vectors. The change of basis to their sum and difference has matrix
\[
 S=\begin{pmatrix}1&1\\1&-1\end{pmatrix},\qquad
 \det S=-2\in O^\times,\qquad
 S^{-1}\begin{pmatrix}0&1\\1&0\end{pmatrix}S=\operatorname{diag}(1,-1).
\]
The new eigenbasis is therefore integral.
In this basis, on \(G_{K'}\) the matrices have the form
\[
                    \begin{pmatrix}u&v\\v&u\end{pmatrix},
       \qquad u^2-v^2=1.
\]
They preserve \(\operatorname{diag}(1,-1)\), and the
other coset acts by multiplying a matrix of this form by
\(\operatorname{diag}(1,-1)\).

In the residue field
the matrices are \(\operatorname{diag}(1,\lambda_0)\)
on the whole rational group. At a place above \(\ell\),
the two eigenvectors are \((1,s)\) and \((1,-s)\)
with \(s=1\) or \(-1\), chosen so that the first inertia
exponent is \(m\). Under conjugation of the place,
the sign is preserved on \(G_{K'}\) and is exchanged
on its other coset. These choices can all be made
over one finite coefficient DVR.

\end{proof}

\subsubsection{Determinant twist for the other factor}

The other rank-two factor will come from an ordinary form. Its determinant is constrained by the four target characters, which fixes the finite twist before the weight is chosen.

Put
\[
             \phi=\eta\xi^\sigma,\qquad
             \Phi=\epsilon_\ell^{k-1}\phi .
\]
The finite character \(\phi\) is odd and is unramified
at \(\ell\). First take a rational prime \(w\) split in \(F_0K\), with
\(w\equiv1\pmod{\ell\operatorname{cond}(\phi)}\) and trivial
residual diagonal. Such primes follow by Chebotarev at the
identity of the compositum. A character \(\theta\) of order
\(\ell\) modulo \(w\) is available and reduces to one.
It is killed on inertia by adjoining the degree-\(\ell\)
subfield of \(\Q(\zeta_w)\).
After fixing this level, take \(k\) odd with
\(k-1\equiv j\pmod{\ell-1}\).
This congruence has odd solutions because \(\xi(-1)=1\)
forces \(j\) even. The Chinese remainder theorem
also lets us impose \(\ell\nmid k\) and any lower
bound determined by the chosen level and \(m_0\).

\subsubsection{Divisibility of the Eisenstein constant}

A cuspidal eigenform with the desired residual system will be obtained by congruence with an Eisenstein series. The constant term must therefore vanish in the residue field of the chosen coefficient ring.

\begin{proposition}[Small Eisenstein constant at the auxiliary conductor]\label{prop:count-three-small-constant}
For \(\chi=\phi\theta\) at the selected prime \(w\), let \(a_0\) denote the constant of the weight-\(k\) Eisenstein series. In the coefficient DVR \(O\), Proposition~\ref{prop:count-three-cusp-congruence} uses the divisibility 
\[
 a_0=\tfrac12 L(1-k,\chi)\in\mathfrak m_O.
\]
 The calculation
uses the original fixed weight, before the final choice of a deep
tower layer.
\end{proposition}

\begin{proof}
Let \(M=\operatorname{cond}(\phi)\) and
\(\chi=\phi\theta\), primitive of conductor \(N=Mw\).
For the Eisenstein series with positive Fourier
coefficients \(\sum_{d\mid n}\chi(d)d^{k-1}\), the
constant term is \(L(1-k,\chi)/2\).
Denote by \(\chi_0\) the character on the same
modulus induced by \(\phi\), so the two characters
have the same reduction. We check that their
special values are integral and congruent.

For any of the characters \(\gamma\) of modulus \(N\)
or \(M\), grouping the Hurwitz zeta series gives
\[
 -k L(1-k,\gamma)=
  (\ell N_\gamma)^{k-1}
       \sum_{a=1}^{\ell N_\gamma}
          \gamma(a) B_k\bigl(a/(\ell N_\gamma)\bigr),
\]
where \(N_\gamma=N\) or \(M\) is the modulus used
in the grouping and \(B_k(X)\) is the Bernoulli
polynomial. Here \(\gamma(a)\) is taken modulo
\(N_\gamma\), even when \(\ell\mid a\).
Expand the polynomial in Bernoulli numbers and
group by \(a\bmod N_\gamma\).
For \(i\ge1\), the coefficient
\((\ell N_\gamma)^{i-1}B_i\) is \(\ell\)-integral by
the von Staudt--Clausen denominator bound and
the integrality of \(B_1=-1/2\) at odd \(\ell\).

For \(i=0\), the coefficient of \(\gamma(b)\)
is the sum of the \(k\)-th powers of all
\(\ell\) representatives of the class of \(b\),
divided by \(\ell N_\gamma\).
The representatives give each residue modulo \(\ell\) once.
Because \(k\) is odd,
\[
          \sum_{r\in\mathbf F_\ell}r^k=0
\]
by pairing \(r\) with \(-r\). This cancels the remaining
factor \(\ell\) in the denominator.
Dividing by \(k\) is allowed as \(\ell\nmid k\),
and at the common modulus \(N\) each coefficient
in this linear expression is independent of
the character. This proves the congruence
of the two special values and integrality of
the special value for \(\phi\).
Finally,
\[
 L(1-k,\chi_0)
     =(1-\phi(w)w^{k-1})L(1-k,\phi)
\]
has positive \(\ell\)-adic valuation. The first
factor is zero modulo \(\ell\) by the choice of
\(w\). Hence the Eisenstein constant belongs to
the maximal ideal of the coefficient DVR.
The weight congruence and this additional choice
of \(k\) leave the eventual lower bound on \(k\) free.
\hfill\mbox{\cite{washington-cyclo,bernoulli-values}}\label{cor:ordinary-count-iii-seed-1}

\end{proof}

\subsubsection{An integral form with unit constant}
\label{par:count-three-cusp-correction}

Divisibility of one Eisenstein constant is only the first step toward a cusp form. An integral correcting form must change that constant while prescribing its values independently at the other cusps.

\begin{proposition}[Cusp correction with unit constant term]\label{prop:unit-constant-cusp-correction}
Let \(N\ge3\), \(\ell\) be odd and prime to \(N\), and
\(\chi(-1)=(-1)^k\). With the separating weight \(m\) of
\eqref{eq:count-three-correction-degree}, assume \(k\ge m+10\).
For the correcting form \(F\), let \(a_n(F;c)\) denote the expansion coefficients at a cusp \(c\), in its chosen integral coordinates. Then 
\[
 a_n(F;c)\in O\quad\text{at every cusp},\qquad a_0(F;\infty)=1,
\]
 with \(O\) the coefficient DVR. This is
the correction used for the primitive Eisenstein series.
\end{proposition}

\begin{proof}
Let \(N\geq3\) with \(\ell\nmid N\), and let
\(\zeta\) be a primitive \(N\)-th root of unity. We use its complex
value for Fourier expansions and its \(\ell\)-adic value for
integrality. For \(a\in(\Z/N)^\times\) and \(r\geq3\), take the
normalized weight-\(r\) Eisenstein sum over \((c,d)\equiv(0,a)\bmod N\)
with summand \((cz+d)^{-r}\). Denote it by \(E_{r,a}\). The normalization is
\[
 E_{r,a}(z)=\frac{N^r(r-1)!}{(-2\pi i)^r}
          \sum_{(c,d)\equiv(0,a)\bmod N}(cz+d)^{-r}.
\]
Absolute convergence gives holomorphy and the right slash action of
\(\gamma=\bigl(\begin{smallmatrix}*&*\\c_\gamma&d_\gamma\end{smallmatrix}\bigr)\)
transports the row \((0,a)\) to \((ac_\gamma,ad_\gamma)\).

In particular the form is modular on \(\Gamma_1(N)\). For
\(\gamma\in\Gamma_0(N)\), the slash is \(E_{r,ad_\gamma}\).
If \(c_\gamma\not\equiv0\bmod N\), the transported sum contains
only terms with \(c\ne0\). Summing the absolute values over \(d\)
gives a bound \(O_r((|c|\operatorname{Im} z)^{1-r})\) for
\(\operatorname{Im} z\geq1\). Writing \(y=\operatorname{Im}z\ge1\) gives the explicit tail bound
\[
 \sum_{c\ne0}(|c|y)^{1-r}
 =2y^{1-r}\sum_{c\ge1}c^{1-r}\longrightarrow0
 \qquad(y\to\infty),
\]
since \(r>2\). The same bound on the
nonconstant part and the convergent \(c=0\) sum give holomorphy
at every cusp.

Set \(L_r(x)=(x\,d/dx)^{r-1}(x/(1-x))\). Differentiating the
partial fraction expansion of \(\pi\cot(\pi t)\) gives
\[
  \frac{(r-1)!}{(-2\pi i)^r}
      \sum_{b\in\Z}(t+b)^{-r}
     =L_r(e^{2\pi i t}).
\]
For \(\operatorname{Im}t>0\) the right side has expansion
\(\sum_{e\geq1}e^{r-1}e^{2\pi i et}\). For the constant
term take \(t=a/N\). For \(c>0\) use \(t=(cz+a)/N\).
The terms with \(c<0\) give the contribution with \(-a\)
and the factor \((-1)^r\). Thus at infinity
\[
 a_0(E_{r,a})=L_r(\zeta^a),\qquad
 a_n(E_{r,a})=\sum_{e\mid n}e^{r-1}
              \bigl(\zeta^{ae}+(-1)^r\zeta^{-ae}\bigr)
       \quad(n\geq1).
\]
The numerator of \(L_r\) is integral and its denominator is
\((1-x)^r\). Reduction is injective on the \(N\)-th roots of
unity, so all these coefficients are integral at \(\ell\).

We next separate the constant at \(a=1\) from that at any
\(a=d\ne\pm1\), with the same weight in the two evaluations.
In the residue field put \(t=x/(1-x)^2\) and
\(u=(1+x)/(1-x)\) for \(x=\bar\zeta^a\). Direct differentiation gives
\[
 L_3(x)=tu,\qquad L_4(x)=t(1+6t),\qquad
 L_5(x)=tu(1+12t),\qquad u^2=1+4t.
\]
At any unit label \(L_3(x)\ne0\), because \(x\ne1,-1\).

In residue characteristic other than two or three, equality of
both \(L_3^2\) and \(L_3L_5\) at \(x,y\) forces
\(t(x)=t(y)\) by taking their quotient.
In characteristic three \(L_4=t\) suffices. Clearing denominators
in \(t(x)=t(y)\) gives \((x-y)(1-xy)=0\).
For \(x=\bar\zeta\), \(y=\bar\zeta^d\), both alternatives
are excluded. At least one of the three families
\[
      E_{4,a},\qquad E_{3,a}^2,\qquad E_{3,a}E_{5,a}
\]
therefore has constants at \(a=1,d\) differing by a unit.

For each such \(d\) choose one family \(H_{d,a}\) of weight \(r_d\).
For a level-one form \(L_{r_d}\) of weight \(r_d\)
with integral coefficients and constant one, set
\[
        H'_{d,a}=H_{d,a}-a_0(H_{d,d})L_{r_d}.
\] Such
forms exist for all even nonnegative weights except two,
by products of \(E_4\) and \(E_6\), including the empty product.
The resulting family has zero constant at \(a=d\) and unit
constant at \(a=1\). Multiplication of the correcting factors gives the even weight
\begin{equation}\label{eq:count-three-correction-degree}
                 m=\sum_{d\ne\pm1}r_d\le8\varphi(N).
\end{equation}
It has zero constants at all \(a\ne\pm1\) and retains a
unit at \(a=1\). The same right slash rule holds for all the
families, since the coefficients in these operations are fixed
as \(a\) varies.

For \(k\geq m+10\), multiply this family by \(E_{3,a}\) if
\(k\) is odd and by \(E_{3,a}^2\) if \(k\) is even. Complete the weight with a product of \(E_4\) and \(E_6\) having constant term one.
The forms \(F_a\) of weight \(k\) so obtained vanish at all
cusps outside the \(\Gamma_0(N)\)-orbit of infinity by the
weight-three factor. At infinity \(c_1=a_0(F_1)\) is a unit. The constant terms
are
\[
       a_0(F_a)=
       \begin{cases}
       c_1,&a=1,\\
       (-1)^kc_1,&a=-1,\\
       0,&a\ne\pm1 .
       \end{cases}
\]
The second case is obtained by slashing by \(-I\).

For \(\chi(-1)=(-1)^k\), set
\[
    G=\frac{1}{2c_1}\sum_{a\in(\Z/N)^\times}\chi(a)^{-1}F_a .
\]
The right slash rule and reindexing \(a\mapsto ad_\gamma\)
give \(G|_k\gamma=\chi(d_\gamma)G\) on \(\Gamma_0(N)\).
Its coefficients are integral because \(2c_1\) is a unit.
Its constant is one, and it vanishes away from the orbit of
infinity. In particular the construction works at the level
\(N=Mw\) with \(w\equiv1\bmod\ell\), even though
\(\ell\mid\varphi(N)\).
\label{cor:ordinary-count-iii-eigenform-1}

\end{proof}

\subsubsection{The cuspidal congruence}

The integral correction can now be multiplied by the small Eisenstein constant. Subtracting it enforces cuspidality while preserving the positive Fourier coefficients modulo the maximal ideal.

\begin{proposition}[A cuspidal congruence at the chosen conductor]\label{prop:count-three-cusp-congruence}
For the primitive Eisenstein character \(\chi\) of conductor \(N=Mw\)
and the small constant calculated for the seed, the cusp correction gives an integral cusp form \(f\) satisfying 
\[
 a_n(f)\equiv\sum_{e\mid n}\chi(e)e^{k-1}\pmod{\mathfrak m}\qquad(n\ge1).
\]
 The
congruence includes the ordinary Hecke eigenvalue at \(\ell\)
used in the lift.
\end{proposition}

\begin{proof}
For the primitive character \(\chi=\phi\theta\) of conductor
\(N=Mw\), the weight-\(k\) Eisenstein series \(E\) of character
\(\chi\) is the sum of \(E_{k,a}\) with weights
\(\chi(a)^{-1}/(2\mathcal G(\chi^{-1}))\) over the unit labels.
Here \(\mathcal G(\chi^{-1})=\sum_a\chi(a)^{-1}\zeta^a\)
is the nonzero primitive Gauss sum, with characters extended
by zero on nonunits. Primitivity gives
\(\sum_a\chi(a)^{-1}\zeta^{ae}=\chi(e)\mathcal G(\chi^{-1})\)
also when \(e\) is a nonunit.

For a nonunit, multiplying
\(a\) by an element fixing \(e\) and on which \(\chi\) is
nontrivial gives zero by primitivity. Parity then gives
\[
    a_n(E)=\sigma(n):=\sum_{e\mid n}\chi(e)e^{k-1}
    \quad(n\geq1),\qquad a_0(E)=L(1-k,\chi)/2 .
\]
For the constant, the same Gauss-sum identity applied
to \(\sum_{e\geq1} e^{k-1}\zeta^{ae}x^e\) for \(|x|<1\)
gives the finite sum of \(L_k(\zeta^a x)\).

At \(x=e^t\), the generalized Bernoulli generating
function gives the continued value
\[
 \left.\left(\frac{d}{dt}\right)^{k-1}
            \frac{\sum_{e=1}^N\chi(e)e^{et}}{1-e^{Nt}}
       \right|_{t=0}=L(1-k,\chi).
\]
The primitive Gauss factor and parity give the normalization of
the constant term above. The cusp vanishing and the character follow
from the corresponding transformations of \(E_{k,a}\).
Set \(f=E-a_0(E)G\). Its constant term at infinity is zero,
and the other positive coefficients satisfy
\[
      a_n(f)\equiv a_n(E)=\sigma(n)\pmod{\mathfrak m}\qquad(n\ge1)
\]
because \(a_0(E)\in\mathfrak m\) and \(G\) has integral coefficients. At cusps in the orbit of infinity,
the vanishing follows by applying the common character law
to the zero constant of \(f\) at infinity. Here the smallness
of \(a_0(E)\) is the earlier Bernoulli and Euler-factor
calculation for the chosen \(k\).

For a prime \(p\), the Hecke operator in this character space has
coefficient \(a_{pn}+\chi(p)p^{k-1}a_{n/p}\), where the second
term is zero if \(p\nmid n\), and \(\chi(p)=0\) if \(p\mid N\).
The divisor sums satisfy the same formula with eigenvalue
\(1+\chi(p)p^{k-1}\). In fact, if \(n=p^b n'\) with
\(p\nmid n'\) and \(v=\chi(p)p^{k-1}\), this is the polynomial
identity
\[
  (1+v+\cdots+v^{b+1})+v(1+v+\cdots+v^{b-1})
      =(1+v)(1+v+\cdots+v^b),
\]
with the second sum empty at \(b=0\).
Thus the reduction of \(f\) is a simultaneous Hecke
eigenvector, nonzero since \(a_1(f)\equiv1\).

\end{proof}

\subsubsection{Simultaneous lifting and the ordinary line}

A nonzero mod-\(\ell\) cusp form gives the desired residual eigenvalues. The second tensor factor requires one characteristic-zero eigenform realizing them simultaneously, including the ordinary eigenvalue at \(\ell\).

\begin{proposition}[Eigensystem of the corrected cusp form]\label{prop:count-three-cuspidal-eigensystem}
The cuspidal congruence of Proposition~\ref{prop:count-three-cusp-congruence}
gives the simultaneous cuspidal eigenform \(g\) and the ordinary local line used for the rank-two factor. After a finite coefficient extension the normalized eigenform satisfies 
\[
 a_p(g)\equiv1+\chi(p)p^{k-1}\pmod{\varpi}\qquad(p\text{ prime}).
\]
 The weight can be
fixed before the choice of the final deep layer.
\end{proposition}

\begin{proof}
\proofstep{An ordinary eigensystem with the selected residual line.}

Choose a finite local coefficient field \(E_\ell\) containing the
coefficient fields of the forms in the construction, and let \(O\)
be its valuation ring with uniformizer \(\varpi\). The
character-\(\chi\) cusp space of level \(N\) has a model over
this field. Indeed normalized newforms have algebraic
coefficients in number fields, and the newform decomposition
spans the cusp space. Since the character has conductor \(N\),
every constituent is new at level \(N\). Both the character
and the coefficients of \(f\) are in the chosen coefficient
field, so \(f\) belongs to this model. One can check this by
taking finitely many Fourier coefficients that identify the
coordinates in a newform basis.

Let \(V_O\) consist of the forms with all coefficients in \(O\).
Take a finite list \(n_1,\ldots,n_s\) separating the forms. Its coefficient map gives
\[
 V_O\hookrightarrow O^s,\qquad h\longmapsto
 (a_{n_1}(h),\ldots,a_{n_s}(h)).
\]
A submodule of a finite free module over a DVR is finite free.
It also contains a basis after integral scaling, by
integrality of the normalized newforms. It is finite free
over \(O\). The coefficient formula makes it stable under
all prime Hecke operators, including those at \(p\mid N\). The coefficient condition gives the saturated reduction
\[
       \ker(V_O\longrightarrow \text{residue expansions})
                                      =\varpi V_O.
\] The mod-\(\varpi\) eigenline of \(f\) therefore
defines a character of the commutative \(O\)-algebra of
Hecke operators on \(V_O\). Write \(\mathfrak m\) for its
kernel. This algebra is finite and torsion-free over \(O\).
Write \(\mathbb T\) for the Hecke algebra. A minimal
prime \(\mathfrak p\subset\mathfrak m\) avoids the uniformizer
since it is a nonzerodivisor.

Normalize \(\mathbb T/\mathfrak p\)
in its fraction field and choose a place over \(\mathfrak m\);
this gives one characteristic-zero evaluation
\[
 \mathbb T\longrightarrow
      \operatorname{Frac}(\mathbb T/\mathfrak p)
\]
with all the prescribed residual Hecke values.
In the generic fibre the faithful module has a nonzero
factor at \(\mathfrak p\). After extending scalars, the
maximal ideal of that Artinian local factor acts nilpotently.
If \(M_\mathfrak p\) is this local generic factor and
\(\mathfrak n\) its maximal ideal, choose the last nonzero
\(\mathfrak n^a M_\mathfrak p\). It is killed by \(\mathfrak n\),
so a nonzero vector there is a common eigenvector.
The eigenvector has nonzero first coefficient, since the
Hecke recursion from a zero first coefficient would give
zero coefficients at every positive index. Normalizing it
gives a cuspidal eigenform \(g\) of level \(N\) and character
\(\chi\) with
\[
    a_p(g)\equiv1+\chi(p)p^{k-1}\pmod{\varpi}
       \qquad(p\text{ prime}),
\]
where \(\varpi\) now denotes a uniformizer after coefficient
extension. This is the simultaneous lifting argument of
Deligne and Serre. \cite[Lemma 6.11]{deligne-serre}

Use the Galois representation of the weight-\(k\) newform
\(g\), with arithmetic Frobenius polynomial
\(X^2-a_p(g)X+\chi(p)p^{k-1}\) for \(p\nmid\ell N\).
It is absolutely irreducible, with determinant
\(D=\Phi\theta\). Since \(\theta\) has order \(\ell\), its residual character is trivial. At every good prime the trace congruence and determinant now give
\[
 \bar P_p(X)=X^2-(1+\bar\Phi(\operatorname{Frob}_p))X+
 \bar\Phi(\operatorname{Frob}_p)
 =(X-1)(X-\bar\Phi(\operatorname{Frob}_p)).
\]
Chebotarev and semisimplification yield \(\bar\Phi\oplus1\). For \(p\mid N\) the local component has conductor
exponent equal to that of the central character, by primitivity
of \(\chi\).

It is principal series with one unramified factor.
Indeed, for a supercuspidal parameter the two-dimensional
irreducible Weil representation has invariants of dimension
zero or two on each normal ramification group. Its determinant
has conductor at most half as large by the Artin conductor
formula. For a special representation the conductor is one
with unramified central character if the twist is unramified.
For a ramified twist of conductor exponent \(a>0\) the
conductor is \(2a\), while that of its central character is
at most \(a\). For two principal series characters with conductors
\(a_1,a_2\), the conductor and determinant bounds read
\[
                a(W)=a_1+a_2,\qquad a(\det W)\le\max(a_1,a_2).
\]
At a ramified prime equality thus forces precisely one
\(a_i\) to be zero. Local--global
compatibility gives the two inertia characters \(1,\phi\theta\).
Outside \(\ell N\) the representation is unramified.
\cite{deligne-galois,carayol-local}

Since \(k\geq3\) and \(\ell\nmid N\), the congruence at
\(p=\ell\) gives \(a_\ell(g)\equiv1\). The unit root
\(\alpha_\ell\) of the Hecke polynomial at \(\ell\) is
congruent to one. The ordinary local theorem for this
newform, or for its unit-root \(\ell\)-stabilization, gives
an unramified quotient \(\nu\) with
\(\nu(\operatorname{Frob}_\ell)=\alpha_\ell\) and invariant
line of character \(D\nu^{-1}\). Its inertia character is
\(\epsilon_\ell^{k-1}\). The two generators of its local data satisfy
\[
         \bar\nu|_{I_\ell}=1,\qquad
         \bar\nu(\operatorname{Frob}_\ell)=1 .
\]
Hence \(\bar\nu\) is trivial on the decomposition group. The high-weight line thus reduces to \(\bar\Phi\)
on that entire group, including when
\(k-1\equiv0\bmod\ell-1\).
For the dual of the attached representation the
zero-exponent line has residual character one and determinant
\(D^{-1}\); the twist by \(D\) gives exactly the preceding
order of the characters. \cite{hida-ordinary-galois}\label{cor:ordinary-count-iii-eigenform-2}

The weight was chosen after \(N\) and the initial positive
exponent \(m_0\) of the quadratic factor. In addition to
the cusp-correction cutoff \(m+10\), we may impose
\(k\geq18m_0+3\) and the earlier congruence and coprimality
conditions. No bound on the eventual deep layer is used in
this choice.
\end{proof}


\subsubsection{The integral lattice and its ordinary line}

The lift and twist fix the rational determinant. The residual endpoint depends also on the lattice in this representation, in particular on which ordinary line survives in its reduction.

\begin{proposition}[The ordinary lattice of the rank-two factor]\label{prop:count-three-ordinary-lattice}
The cuspidal lift and determinant twist give an integral rank-two representation \(r\) with 
\[
 \det r(g)=\Phi(g)\qquad(g\in\ker\theta).
\]

One can choose its lattice and the marked eigenbasis so that
the ordinary local line has the reduction and ordered
character required for the tensor.
\end{proposition}

\begin{proof}
\proofstep{Choose the ordinary integral lattice.}

The integral rank-two datum we now construct is
\(\rho_1\) with the following properties. Its determinant
is \(\Phi\) on \(\ker\theta\). At \(g_0\) it is diagonal,
and its reduction on \(G_{\Q}\) is
\[
                  \begin{pmatrix}\bar\Phi&0\\b&1\end{pmatrix}
\]
with \(b\ne0\). Away from \(\ell\) it is the identity on
inertia where both \(\phi\) and \(\theta\) are trivial.
At each chosen place \(A\) above \(\ell\), it has an
integral invariant line generated by \((1,x_A)\),
with ordered inertia characters \(\epsilon_\ell^{k-1},1\).
Start with the cuspidal attachment just constructed. Its
ordinary invariant line is the character \(D\nu^{-1}\), with
residue \(\bar\Phi\) on the entire decomposition group.

Choose a stable lattice and use the two eigenprojectors
of \(\rho(g_0)\). If the two eigenvalues are \(\lambda_1,\lambda_2\), the projectors are
\[
 e_i=\frac{\rho(g_0)-\lambda_j I}{\lambda_i-\lambda_j}
 \qquad(\{i,j\}=\{1,2\}).
\]
The denominator is a unit by residual separation, so both preserve the lattice. Order the basis
with the \(\bar\Phi(g_0)\)-eigenvalue first. The traces
of \(\rho(g)\) and of \(\rho(g_0g)\) determine its
two reduced diagonal entries as \(\bar\Phi(g),1\).
The first diagonal entry in a product is
\(a_ga_h+b_gc_h\) if \(b_g,c_h\) are the off-diagonal
entries. Comparing with the reduced diagonal character gives
\[
                  \overline{b_gc_h}=0\qquad(g,h\in G_\Q).
\]

The two sets of off-diagonal entries have finite
minimum valuations \(v_+,v_-\), by irreducibility,
and \(v_++v_->0\). Multiplying the second basis vector by \(\varpi^{v_-}\)
changes the minima to
\[
           v'_-=v_--v_-=0,\qquad v'_+=v_++v_->0.
\]
Thus a lower-left entry survives and the upper-right ones all
vanish after reduction.
At each place over \(\ell\), a rational Frobenius
lift with trivial cyclotomic value separates \(\bar\Phi\)
from one, by Paragraph~\ref{par:dyadic-residual}.
The reduced ordinary line is the first, \(\bar\Phi\)-eigenline there. A generator \((u,v)\) of the lifted line has \(u\) a unit and admits the normalization
\[
 u^{-1}(u,v)=(1,x_A),\qquad x_A=v/u.
\]


The cuspidal rank-two representation is irreducible
on every open subgroup. Indeed, a reducible
restriction to an open normal subgroup that contains
a power of an inertia element with distinct eigenvalues
has at most two invariant lines. Such an element
exists here since the local cyclotomic exponents are
distinct. By irreducibility on the rational group
there must be two lines, exchanged by the other
coset of their stabilizer. On the stabilizer,
the two reduced line characters are \(1,\bar\Phi\).
Conjugation exchanges them and fixes \(\bar\Phi\),
so \(\bar\Phi\) is one on the stabilizer. Writing \(H\) for the index-two stabilizer, each square belongs to \(H\). Thus
\[
 \bar\Phi(g)^2=\bar\Phi(g^2)=1\qquad(g\in G_\Q).
\] At the mark this contradicts
\(\bar\Phi(g_0)^2=z^2\ne1\). The latter follows from
\(\bar\Phi(g_0)=-z\), with \(z\) of nontrivial odd
order. The argument for arbitrary open subgroups
follows by passing to the normal core.

Put \(u=(\epsilon_\ell/\omega)^{1/2}\), the root with
values in \(1+\ell\Z_\ell\), and define
\[
 \delta =
 \begin{cases}
 u^{j+2-k} & \sigma=1,\\
 \omega^j\xi\,u^{j+2-k} & \sigma=-1 .
 \end{cases}
\]
The identities \(\omega^{k-1}=\omega^j\) and
\(\omega^{2j}=1\) give
\(\delta^2\Phi=\vartheta\).
For the original order, \(\bar\delta=1=d_3\) and
\(\bar\delta\bar\Phi=d_0=\bar\vartheta\).

For the exchanged order, \(\bar\delta=\bar\Psi=d_3\)
and \(\bar\delta\bar\Phi=d_0=\bar\alpha\).
Together with \(d_0/d_1=\lambda_0\) and
\(d_0d_3=d_1d_2\), this identifies all four reductions.
On any inertia group away from \(\ell\) where the
residual diagonal is trivial, \(\omega=1\) and the
values of both finite characters \(\eta,\xi\) are
one by prime-to-\(\ell\) injectivity. Thus
\(\delta=\phi=1\) there.
The same conclusions hold on the part of inertia
above \(\ell\) in the cyclotomic kernel, since
\(\eta,\xi\) are unramified above \(\ell\).

\end{proof}

\subsubsection{Tensor representation and residual endpoint}

We can now tensor the two integral bases in the marked order. This realizes the endpoint as an actual residual representation with a flag, data that a list of four characters alone would not specify.

\begin{proposition}[Integral tensor and its marked endpoint]\label{prop:count-three-tensor-endpoint}
With the ordinary bases of the two factors, their twisted tensor
is symplectic and has the required ordered flag on the fitted
field. Its reduction gives the designated endpoint and the local diagonal characters of the rank-four chain. Write \(b_i\) for the inertia exponent on the \(i\)th ordered line after the scalar twist. Then 
\[
 2(b_1,b_2,b_3,b_4)=2(k-1+m,k-1-m,m,-m)+(j+2-k)(1,1,1,1).
\]

\end{proposition}

\begin{proof}
Let \(E\) be the free rank-two module of \(\rho_1\), with
the standard alternating form, and \(H\) that of \(\tau\)
with symmetric form \(\operatorname{diag}(1,-1)\).
For bases \((e_0,e_1)\), \((f_0,f_1)\), take the tensor
basis
\[
               (e_0 f_0,e_0 f_1,-e_1 f_1,e_1 f_0).
\]
The tensor form is alternating and has the prescribed
pairings in indices \(0,3\) and \(1,2\). In this basis,
\(\delta(\rho_1\otimes\tau)\) has similitude
\(\delta^2\det\rho_1\). Thus after killing \(\theta\)
on the deep layer its similitude is exactly
\(\vartheta\). On the tame layer the same equality
holds after reduction, since \(\theta\) reduces to one.
Both factors are diagonal at the normalizer and all
its powers.

\proofstep{Compute the tensor flag in the ordered basis.}

For the local flag let \(L_x=\begin{pmatrix}1&0\\x&1\end{pmatrix}\).
Conjugate the tensor first by \(L_{x_A}\otimes1\) in
the displayed basis. Then conjugate by
\[
             1+s_A E_{10}-s_A E_{32},
\]
where \(s_A\) is the sign of the chosen eigenline of
\(\tau\) at \(A\). Both are lower unipotent symplectic
matrices. Their product is the flag matrix, with
\((1,0)\) entry \(s_A\).
Over characteristic zero the local conjugate is upper
triangular, since the first step triangularizes the
rank-two factor and the second orders the two lines
of the symmetric factor in each block. Its inertia
characters, in order, are
\[
       \delta\epsilon_\ell^{k-1+m},\quad
       \delta\epsilon_\ell^{k-1-m},\quad
       \delta\epsilon_\ell^m,\quad
       \delta\epsilon_\ell^{-m}.
\]

On the cyclotomic kernel the characters are all one.
They reduce to the required characters on the full
marked decomposition groups. To see the last point,
the marked place splits in \(K'\), so \(d_0=d_1\)
and \(d_2=d_3\) there. Conjugation by \(L_{x_A}\) preserves the residual lower
triangular order. Invariance of the first line then forces its
lower-left entry to vanish, so its reduction remains
\[
                          \operatorname{diag}(\bar\Phi,1).
\]
This proves the local diagonal conditions.

Outside \(\ell,w\), the residual tensor kills the
required inertia. Indeed where the original
residual diagonal is trivial, \(\lambda_0=1\) on
the restricted group. There \(\tau\) is unramified,
and \(\phi=\delta=\theta=1\) on the indicated inertia.
On the deep layer the corresponding integral claim
follows because the extra inertia has been killed.
The construction gives the tame scout \(S\) and
the compatible integral matrices on every such
layer, since they all come by restriction from the
same rational tensor. The chosen frames over
different layers transport as the lower factors
for this same ordinary flag. In the quadratic
factor the place sign changes by \(\lambda_0\)
on conjugation, and the rank-two line transports
by \(\rho_1\). These two actions give the flag
transport simultaneously on both factors.

In the residue field, the matrix is
\[
       \bar\delta\,
       \begin{pmatrix}
        \bar\Phi&0&0&0\\
        0&\lambda_0\bar\Phi&0&0\\
        0&-\lambda_0 b&\lambda_0&0\\
        b&0&0&1
       \end{pmatrix}.
\]
Thus only \(-w_\beta,-w_\vartheta\) can occur
globally off the diagonal. The lower frame entry
\(s_A\) stays nonzero. If \(b\) became zero on
\(G_F\), then conjugating the powered normalizer
by any rational Galois element would give
\[
           0=b(g)\bigl(\bar\Phi(g_0^{\ell^r})-1\bigr).
\]
This uses the product rule
\(b(gh)=b(g)\bar\Phi(h)+b(h)\), and the powered
normalizer still lies in \(G_F\) with
\(\bar\Phi\ne1\). Indeed \(\bar\Phi\) at the normalizer
is the ratio \(d_0/d_3\). Hence \(b\) is nonzero
on \(G_F\), and both indicated global roots occur.
The tensor supplies exactly the sign and support
conditions of Section~\ref{sec:count-three-chain}.

Finally, the doubled inertia exponents of this
tensor are the four integers
\[
 (2(k-1+m),2(k-1-m),2m,-2m)+(j+2-k)(1,1,1,1).
\]
Their adjacent differences are
\(4m,2(k-1-2m),4m\).
They are positive even integers at least four if
\(k\ge2m+3\). Removing the scalar by twisting
by \(\delta^{-1}\) changes the doubled inertial
exponent by \(-(j+2-k)\). Its sum with the first
doubled weight is even, as needed for the
ordinary seed comparison.
\hfill\mbox{\cite{ray-class}}\label{cor:ordinary-count-iii-seed-2}

\end{proof}

\subsubsection{Irreducibility after cyclotomic restriction}

To obtain an automorphic starting point we still need irreducibility after cyclotomic restriction and regularity of the weights after opposite shifts of the quadratic factor.

\begin{proposition}[Irreducibility and separation of infinity exponents]\label{prop:count-three-tensor-irreducible}
If the rank-two cuspidal factor is absolutely irreducible on all
open subgroups, the tensor restricted to the fitted CM layer
remains absolutely irreducible after killing the cyclotomic
character. For \(k\ge18m_0+3\) there is \(1\le a\le9\) for which the
opposite shifts \(\pm am_0\) give four distinct infinity parameters
and 
\[
 k-1-2am_0\ge2.
\]

\end{proposition}

\begin{proof}
Suppose the rank-two cuspidal input is absolutely
irreducible on all open subgroups of \(G_{\Q}\).
Put \(L=FK\) and \(M=LK'\).
Let \(H=G_M\cap\ker\epsilon_\ell\). If \(H\) contains a
nonscalar matrix, an invariant line for \(H\) belongs to the
at most two eigendirections of that matrix. For an invariant line \(l\), normality of \(H\) makes every \(G_M\)-translate an \(H\)-invariant line. Hence
\[
 [G_M:\operatorname{Stab}(l)]=|G_M l|\le2 .
\]
Continuity makes the stabilizer closed, so this finite-index subgroup is open and preserves \(l\).
If the subgroup acts by scalars, the projective
action factors through the procyclic compact
cyclotomic image of \(G_M\). A lift of a
topological generator preserves an eigenline,
and so does the closure of its projective cyclic
group. In both cases this contradicts the assumed
strong irreducibility.

The character of \(M/L\) is nontrivial on
\(G_L\cap\ker\epsilon_\ell\). Otherwise this quadratic
character would factor through \(\epsilon_\ell(G_L)\)
and kill its pro-\(\ell\) part. The powered normalizer
lies in \(G_L\) with \(\omega=1\) and
\[
 \lambda_0=-1,
\]
a contradiction.
Over \(G_M\cap\ker\epsilon_\ell\), the tensor without
its scalar is the sum of the two irreducibles
\(\rho\otimes(\kappa')^a\) and
\(\rho\otimes(\kappa')^{-a}\).
They are inequivalent. Choose an inertia element
\(i\in G_M\) above \(\ell\) with infinite cyclotomic
order, and let \(g\in G_L\) be the powered
normalizer outside \(G_M\). The commutator belongs to \(G_M\) by normality and to \(\ker\epsilon_\ell\) by abelianity. Since conjugation by \(g\) inverts \(\kappa'\), its value is
\[
 \kappa'(igi^{-1}g^{-1})
 =\kappa'(i)\kappa'(gi^{-1}g^{-1})
 =\kappa'(i)^2 ,
\]
of infinite order. On this commutator the ratio of the determinants of
the two rank-two summands is
\[
 \frac{\det(\rho\otimes(\kappa')^a)}
           {\det(\rho\otimes(\kappa')^{-a})}
       =(\kappa')^{4a},
\]
whose value there is
\[
 \kappa'(i)^{8a}\ne1.
\]

Thus the summands are inequivalent. The restriction to
\(G_M\cap\ker\epsilon_\ell\) is a direct sum with two
inequivalent simple summands, so a nonzero proper
invariant subspace for that subgroup must be one summand.
The other coset in \(G_L\cap\ker\epsilon_\ell\) exchanges
them, so neither is invariant for the whole group.
This also proves irreducibility on \(G_L\).
Scalar twisting has no effect on these conclusions.

For the infinity exponents in the induction
argument, the two conjugate places of \(K'\)
have two rank-two parameters each. Twisting
by \((\kappa')^a\) shifts these two parameter sets
by opposite integers \(am_0\) and \(-am_0\).
Only differences of the original parameters
in the two sets can equal \(2am_0\).
Writing the unshifted entries as \(b_1,b_2\) and \(c_1,c_2\), an overlap in either orientation requires
\[
  2am_0\in\{\pm(c_j-b_i):1\le i,j\le2\}.
\]
This set has at most eight elements, whereas \(2am_0\) takes nine distinct values in the indicated range. One of \(a=1,\ldots,9\)
therefore separates the four parameters. For the adjacent gaps,
\[
                       k-1-2am_0\ge2
\]
follows from the earlier choice \(k\ge18m_0+3\).
\label{cor:ordinary-count-iii-seed-3}

\end{proof}

\subsubsection{Cuspidal induction on the fitted CM layer}
\label{par:count-three-cuspidal-induction}

The separation of the infinity exponents is used here to obtain a regular cuspidal representation. Its restriction to the fitted CM field supplies the automorphic realization of the integral tensor.

\begin{proposition}[Cuspidal induction of the separated tensor]\label{prop:count-three-cuspidal-induction}
For the chosen infinity separation and the tensor of
Propositions~\ref{prop:count-three-tensor-endpoint}
and \ref{prop:count-three-tensor-irreducible}, the restriction to the fitted CM layer \(L\) has the regular algebraic polarized cuspidal attachment used for the seed. For the separated tensor before the common scalar twist, the ordered infinity labels and their consecutive gaps are 
\[
 (k-1+m,k-1-m,m,-m),\qquad(2m,k-1-2m,2m).
\]

\end{proposition}

\begin{proof}
With \(a\) chosen by the infinity separation, put
\(L=FK\) and \(M'=LK'\). The latter extension has degree
two over \(L\) because the powered normalizer belongs
to \(G_L\) and acts nontrivially on \(K'\). The extension
\(M'/K'\) is solvable, since both \(L\) and \(K'\) are
abelian over \(\Q\). Let \(\pi\) be the cuspidal representation
of \(\mathrm{GL}_2(\mathbb A_\Q)\) of the newform \(g\).
Base change to \(K'\), twist by the Hecke character of
\((\kappa')^a\), and base change to \(M'\). These are
cuspidal base changes.

At each cyclic step, failure of
cuspidality would give a nontrivial self-twist by a character
of the cyclic extension. Comparison at unramified places would give a Galois self-twist by a nontrivial quotient character \(\nu\). An intertwiner \(T\) restricts on the upper subgroup \(H\) to an element of its scalar commutant. Hence
\[
 T=cI,\qquad T r(g)T^{-1}=\nu(g)r(g)
 \quad\Longrightarrow\quad \nu(g)=1\quad\text{for all }g,
\]
which is impossible. This holds by the rank-two strong irreducibility
and is preserved by the character twist. At unramified places, base change takes the powers
dictated by the residue degrees and the twist multiplies
by the rank-one value. Thus the attached representation is
\[
                 \rho|_{G_{M'}}\otimes(\kappa')^a .
\]
\cite{clozel-rajan}\label{cor:ordinary-count-iii-admission-1}

Write \(R=\rho\otimes\operatorname{Ind}_{G_{K'}}^{G_\Q}
(\kappa')^a\), with the tensor basis already specified.
The tensor identity for induction gives
\[
  R|_{G_L}
   \simeq\operatorname{Ind}_{G_{M'}}^{G_L}
                 \bigl(\rho|_{G_{M'}}(\kappa')^a\bigr).
\]
The two constituents on \(G_{M'}\) are inequivalent by
the determinant test above. Hence the cuspidal representation
over \(M'\) is inequivalent to its conjugate over \(L\),
by unramified Galois comparison. Quadratic automorphic
induction is cuspidal under this hypothesis. At an embedding \(\tau:L\hookrightarrow\C\), if \(\tau_+,\tau_-\)
are the two extensions to \(M'\), the infinity multiset is
\(\Lambda_{\tau_+}\sqcup\Lambda_{\tau_-}\). If \(B_+,B_-\) denote the unshifted rank-two multisets at the two embeddings, the character of \(K'\) gives
\[
 \Lambda_{\tau_+}\sqcup\Lambda_{\tau_-}
 =(B_++am_0)\sqcup(B_--am_0)
\]
after choosing the positive embedding first.

Their
restrictions to \(K'\) are the two distinct embeddings,
so the parameters are the four separated by the choice
of \(a\). They are algebraic in rank four: the rank-two
algebraic shifts have the same half-integral parity as
the rank-four shifts, and the twists add integers.
Thus the induced representation is regular algebraic,
with \(R|_{G_L}\) as its Galois representation by comparison
at the unramified places. For the polarization, write \(W=\rho|_{G_{M'}}\otimes(\kappa')^a\) and let \(\iota\) denote complex conjugation. The rank-two factor extends over \(\Q\), its determinant is \(D\), and \((\kappa')^\iota=(\kappa')^{-1}\). The exterior pairing in rank two therefore gives
\[
 W^\iota\simeq \rho|_{G_{M'}}\otimes(\kappa')^{-a}
       \simeq W^\vee\otimes D|_{G_{M'}} .
\]
Induction and comparison at unramified places give the corresponding essential conjugate duality of the automorphic representation by strong multiplicity one.
\cite{clozel-rajan}\label{cor:ordinary-count-iii-admission-2}

The finite-place conditions can be checked before induction.
At primes in \(S\) and at \(w\), the extension \(M'/L\)
splits, by the splitting in \(K'\) at \(S\) and the
choice at \(w\). At primes away from \(S\cup\{w\}\), the
quadratic extension \(M'/L\) is unramified. In fact the
quadratic character there is the restriction of \(\bar\alpha\)
on \(G_L\), and
\[
 \alpha=\eta\epsilon_\ell/\omega
\]
is
unramified away from \(S\). The representation
\(\rho|_{G_L}\) is unramified at these places as well, since
the conductors of \(\phi\) and \(\theta\) are supported on
\(S\cup\{w\}\). The Hecke character of \(K'\) is
unramified away from \(\ell\), so the twisted base change
is spherical at the nonsplit places of \(M'/L\).
Here one uses the local compatibility of the rank-two
attachment away from \(\ell\), under restriction and
character twist.

At a split place of \(M'/L\) the
parameter of the induction is the direct sum from its
two places. At a nonsplit unramified place the inertia group \(I\) is the same for the quadratic extension and the base. On restricting the induced parameter to inertia, its two coset summands both restrict trivially:
\[
 \left.(\operatorname{Ind}\tau)\right|_I
 =\left.\tau\right|_I\oplus\left.\tau^c\right|_I
 =\mathbf{1}^{\oplus4}.
\]
The summands have zero monodromy, hence the induction is spherical there. These are
the split and unramified local comparisons for the induced representation.

In particular the induced rank-four representation matches
the tensor at all places away from \(\ell\), including
monodromy. At \(\ell\) the quadratic extension is split. The ordered
exponents of the tensor before the scalar shift and their gaps are
\[
      (k-1+m,k-1-m,m,-m),\qquad
      (2m,k-1-2m,2m).
\]
Every gap is at least two. The local Weil lines are unlinked
by the valuation calculation in
Section~\ref{sec:count-two-seed-specialization}.
The local--global compatibility theorem at \(\ell\)
applies to the regular algebraic polarized cuspidal
representation, since those gaps give Shin regularity.
\cite{blggt,blggt-lg}\label{cor:ordinary-count-iii-admission-3}

\end{proof}

\subsubsection{The eigenpoint at both ordinary levels}
\label{par:count-three-actual-seed}

The cuspidal tensor has to give an ordinary specialization at each of two levels. With the chosen local flags one can impose both sets of Hecke relations at the same integral point.

\begin{proposition}[Support of the integral tensor seed]\label{prop:count-three-ordinary-attachment}
The normalized integral tensor gives a seed evaluation on both
ordinary Hecke quotients. In the framed ring over \(F\), write \(f_i\) for their level maps. For these maps \(f_i\) and ideals \(J_i\), every prime \(P\) below the seed kernel in \(B\) satisfies 
\[
 f_1(J_1)B+f_2(J_2)B\subseteq P\subseteq\ker y_{\rm seed}.
\]
 This is the initial support condition for the four
residual links.
\end{proposition}

\begin{proof}
On the framed integral tensor the scalar is \(\delta\);
the twist over \(L\) by \(\delta^{-1}\) recovers
\(\delta^{-1}(\delta R|_{G_L})=R|_{G_L}\), the attached
cuspidal representation. Write \(e_0=j+2-k\)
for the exponent of \(\delta^2\) on an open subgroup.
There is a positive integer \(h\) with
\[
 (\delta^{-2}\epsilon_\ell^{e_0})^h=1,
\]
since this
character has finite image. Along with the even doubled
gaps and the parity of the doubled first exponent plus
the doubled correction, this is the scalar correction data
in the ordinary seed interface. The required irreducibility
on the cyclotomic kernel was proved above. The choice of
\(\delta\) gives determinant multiplier \(\vartheta\) on \(G_F\)
because \(\theta\) is trivial there.

For each link use the parallel weight shift of
Paragraph~\ref{par:count-two-seed-support}, which is
needed here as the seed's doubled weights can be odd.
Write \(\beta_i\) for the link companion and \(h_0\) for
the first doubled seed exponent. Take the shift
\(n=|h_0|\). All adjacent doubled gaps are even. Thus for every \(r\)
\[
 h_r\equiv h_0\equiv |h_0|=n\pmod2,\qquad (h_r+n)/2\in\Z .
\]
This gives integral seed exponents for \(\kappa\beta_i^n\) at the chosen embeddings. Thus on \(G_L\) the shifted
evaluation has the cuspidal attachment just constructed
twisted by the algebraic Hecke character of
\(\kappa\beta_i^n\delta\).

At a selected embedding, if \(m^{(R)}_r\) denotes the
\(r\)-th integral untwisted tensor exponent, the scalar exponent is
\[
             \frac{h_r+n}{2}-m^{(R)}_r\in\Z .
\]
It is independent of \(r\), since the companion acts by one
scalar on all four lines.
At the conjugate embeddings integrality follows from
the link pairing with its integral cyclotomic multiplier
and the rank-two determinant \(D\). The anticyclotomic
companion preserves the link pairing and its unitary
normalization.

For the tame types, after killing \(\theta\) on \(G_F\)
the inertia of \(R\) away from \(\ell\) has finite image of
order prime to \(\ell\). Where \(K'\) splits this follows
from the rank-two inertia calculation and unramifiedness
of \(\kappa'\). At the other places the quadratic extension
is unramified and the whole seed is unramified. The
finite inertia of \(\kappa\beta_i^n\delta\) away from \(\ell\) has
order prime to \(\ell\), by the common twist and the
formula for \(\delta\). If \(e\) is the order of this finite inertia image, each inertial root satisfies
\[
 z^e=1,\qquad e\in\mathbb Z_\ell^\times.
\]
Thus the polynomial \(X^e-1\) has simple roots modulo \(\ell\). Each inertial root \(z\) is consequently determined by its reduction:
\[
 z^e=1,\qquad (X^e-1)'(z)=e z^{e-1}\in O^\times
 \quad\Longrightarrow\quad z=[\bar z].
\]
Here the brackets denote the Teichmüller lift. This is the type at each prescribed
tame place, with zero monodromy.

At \(w\) it is unramified
after killing \(\theta\). At both exceptional places the
seed is unramified as required: these places are outside
\(S\cup\{w\}\), and their inertia is killed by the common
link twist. The finite places where it can ramify are
split over \(F\), including those over \(\ell\).
Above \(\ell\), the same exponent gaps and the ordered
local characters from the tensor flag give the ordinary
vector using the ordered translation in
Section~\ref{sec:count-two-seed-specialization}.
At inert places over \(F\) the attachment is spherical.
The compact unitary descent and the projection to
ordinary scalar functions therefore apply with exactly
the prescribed types and one exceptional place on each
of the two levels.
\cite{thorne-unitary-descent,newton-thorne-unitary}\label{cor:ordinary-count-iii-admission-4}

Let
\[
 f_{i,n}:R_i\to B
\]
be the shifted link map to the
common framed ring and \(y_{\rm seed}\) the tensor point.
The characteristic coefficients of \(y_{\rm seed}\circ f_{i,n}\)
agree with those of the ordinary eigenpoint at split good
primes, hence everywhere by density. The local-character
coordinates agree by the matched weights and uniformizer
translations. The spectral and local matches give
\[
 J_i\subseteq\ker(y_{\rm seed}\circ f_{i,n}).
\]

By the finite-twist invariance in
Paragraph~\ref{par:count-two-seed-support},
\[
                  f_{i,n}(J_i)B=f_i(J_i)B
                             \subseteq\ker y_{\rm seed}.
\]
This transfers support to the original link. The hypotheses of the adjoint comparison in the same
paragraph hold on this tensor: it is cuspidal, irreducible
on the cyclotomic kernel, has the separated ordinary
exponents, and is unramified at the exceptional place.
The comparison gives vanishing at every prime below
the seed evaluation kernel. In the normalized framed
ring over \(F\), the sum of the two extended ideals therefore
satisfies the seed condition for the four links.
Propositions~\ref{prop:tower-residual-propagation},
\ref{prop:count-three-residual-path} and
\ref{prop:count-three-extra-vertices} propagate this same ideal
to the normalized target in the rank-four construction.
\end{proof}

